%! Solving Radical Equations \& Extraneous Solutions — Unit 6, Lesson 6.5 — Group Activity
\documentclass[10pt]{article}
\usepackage{algebra2-article}
\usepackage{algebra2-boxes}
\newcommand{\TallMath}[1]{$\displaystyle #1\rule[-1.4em]{0pt}{3.2em}$}
\newcommand{\lgrid}[4]{%
  \draw[step=1, linegray!40, very thin] (#1,#3) grid (#2,#4);
  \draw[->, semithick, charcoal] (#1,0)--(#2,0) node[below right, font=\scriptsize]{$x$};
  \draw[->, semithick, charcoal] (0,#3)--(0,#4) node[left, font=\scriptsize]{$y$};}

\begin{document}

\pageheader{Unit 6, Lesson 6.5}{Group Activity}
\namepartnerperiod

\vspace{0.08in}

\begin{headlinebox}{forestbg}
\small\textbf{Candidate or Solution?} Every answer you produce today starts life as a
\emph{candidate}. Only the check promotes it. With your partner, solve each equation, check
\emph{every} candidate in the original, and say out loud which ones survived --- and why the others
did not.
\end{headlinebox}

\vspace{0.06in}

\begin{tcolorbox}[colback=white, colframe=black!40, title=\textbf{Tier R --- Remediate},
    fonttitle=\bfseries, arc=1.5mm, left=3mm, right=3mm, top=2mm, bottom=2mm, breakable]
\textbf{Part 1 --- The four steps, one at a time.} \; Solve $\sqrt{x-1}+2=5$.
\begin{center}
\renewcommand{\arraystretch}{1.8}
\begin{tabularx}{0.94\linewidth}{@{}l X@{}}
\textbf{Step 1.} Isolate the radical. & $\sqrt{x-1}=\blank{2.0cm}$ \\
\textbf{Step 2.} Square both sides. & $x-1=\blank{2.0cm}$ \\
\textbf{Step 3.} Solve. & $x=\blank{2.0cm}$ \\
\textbf{Step 4.} Check in the \emph{original}. & $\sqrt{\blank{1.2cm}-1}+2=\blank{1.0cm}+2=\blank{1.0cm}$ \\
\end{tabularx}
\end{center}
Solution set: \blank{2.4cm}

\vspace{5pt}
\textbf{Part 2 --- Now you run the steps.} Solve and check each one.
\begin{enumerate}[leftmargin=1.8em, itemsep=6pt]
  \item $\sqrt{4x+1}=5$ \hfill $x=\blank{1.6cm}$ \quad Check: \blank{3.0cm}
        \vspace{0.30in}
  \item $\sqrt[3]{x+3}=-1$ \hfill $x=\blank{1.6cm}$ \quad Check: \blank{3.0cm}
        \vspace{0.30in}

        A cube root came out \emph{negative} and the answer was still fine. Why is that allowed
        here, when it would not be allowed for a square root? \blank{6.0cm}
  \item $\sqrt{x+2}+7=3$
        \vspace{0.10in}

        Isolate first: $\sqrt{x+2}=\blank{1.4cm}$. \; Stop and look at that line.
        \textbf{Answer:} \blank{3.4cm}

        Explain in one sentence how you knew without doing any more work. \blank{6.0cm}
\end{enumerate}

\vspace{4pt}
\textbf{Part 3 --- Confirm it on the graph.} Below is $y=\sqrt{x-1}+2$ with two horizontal lines.
\begin{center}
\begin{tikzpicture}[scale=0.44]
  \lgrid{-1}{14}{-1}{7}
  \begin{scope}
    \clip (-1,-1) rectangle (14,7);
    \draw[very thick, forest, domain=1:14, samples=110, smooth] plot (\x,{sqrt((\x)-1)+2});
    \draw[very thick, navy] (-1,5)--(14,5);
    \draw[very thick, goldacc] (-1,1)--(14,1);
  \end{scope}
  \fill[forest] (1,2) circle (4pt);
  \fill[charcoal] (10,5) circle (3pt);
  \node[navy, font=\scriptsize, right] at (11.4,5.5) {$y=5$};
  \node[goldacc, font=\scriptsize, right] at (11.4,1.5) {$y=1$};
\end{tikzpicture}
\end{center}
\begin{enumerate}[label=(\alph*), leftmargin=1.9em, itemsep=5pt, start=4]
  \item The curve meets $y=5$ at $x=\blank{1.4cm}$. Does that match your Part 1 answer?
        \blank{1.4cm}
  \item How many times does the curve meet $y=1$? \blank{1.4cm} \; So how many solutions does
        $\sqrt{x-1}+2=1$ have? \blank{1.4cm}
\end{enumerate}
\end{tcolorbox}

\begin{tcolorbox}[colback=white, colframe=black!40, title=\textbf{Tier A --- Approaching Proficiency},
    fonttitle=\bfseries, arc=1.5mm, left=3mm, right=3mm, top=2mm, bottom=2mm, breakable]
\textbf{Part 1 --- Sort the candidates.} Solve each equation. Write \emph{every} candidate the
algebra gives you, check each one, and then state the solution set. Show your squaring work in the
space below the table.
\begin{center}
\renewcommand{\arraystretch}{1.7}
\begin{tabularx}{\linewidth}{@{}l X X X@{}}
\hline
\textbf{Equation} & \textbf{Candidates} & \textbf{Which ones check?} & \textbf{Solution set} \\
\hline
(a) \; $\sqrt{x+9}=x+3$ & \blank{2.2cm} & \blank{2.2cm} & \blank{1.8cm} \\
(b) \; $\sqrt{2x+5}=x+1$ & \blank{2.2cm} & \blank{2.2cm} & \blank{1.8cm} \\
(c) \; $\sqrt{x-3}=x-5$ & \blank{2.2cm} & \blank{2.2cm} & \blank{1.8cm} \\
(d) \; $\sqrt[3]{2x-1}=3$ & \blank{2.2cm} & \blank{2.2cm} & \blank{1.8cm} \\
(e) \; $\sqrt{x+6}+4=1$ & \blank{2.2cm} & \blank{2.2cm} & \blank{1.8cm} \\
\hline
\end{tabularx}
\end{center}
\vspace{1.9in}

\textbf{Part 2 --- Find the pattern.} Look back at the candidates you \emph{rejected} in (a), (b),
and (c).
\begin{enumerate}[label=\textbf{J\arabic*.}, leftmargin=2.4em, itemsep=6pt]
  \item For each rejected candidate, work out the value of the \emph{right-hand side} of the
        original equation.
        \begin{center}
        \renewcommand{\arraystretch}{1.5}
        \begin{tabularx}{0.9\linewidth}{@{}X X X@{}}
        \hline
        \textbf{Rejected candidate} & \textbf{Right side equals} & \textbf{Its sign} \\
        \hline
        (a) $x=\blank{1.4cm}$ & \blank{1.6cm} & \blank{1.6cm} \\
        (b) $x=\blank{1.4cm}$ & \blank{1.6cm} & \blank{1.6cm} \\
        (c) $x=\blank{1.4cm}$ & \blank{1.6cm} & \blank{1.6cm} \\
        \hline
        \end{tabularx}
        \end{center}
        Write the pattern as a rule someone could use: \emph{a candidate is extraneous exactly
        when} \dots \writelines{2}
  \item Equation (d) was a cube root, and it produced no extraneous candidate. Explain why it
        \emph{could not have}, using the words \textbf{index} and \textbf{sign}. \writelines{3}
  \item In (e) you never had to square anything. What did you see, and why is spotting it worth
        more than the two minutes it saves? \writelines{2}
\end{enumerate}
\end{tcolorbox}

\begin{tcolorbox}[colback=white, colframe=black!40, title=\textbf{Tier E --- Extension},
    fonttitle=\bfseries, arc=1.5mm, left=3mm, right=3mm, top=2mm, bottom=2mm, breakable]
\textbf{Part 1 --- Two radicals.} The SOL only asks you to handle one radical. Here is what happens
with two.
\begin{enumerate}[label=(\alph*), leftmargin=1.9em, itemsep=6pt]
  \item $\sqrt{3x+1}=\sqrt{x+9}$ \; --- both sides are already isolated radicals, so square once.
        \vspace{0.45in}

        $x=\blank{1.6cm}$ \quad Check: \blank{4.0cm}
  \item $\sqrt{x+5}-\sqrt{x}=1$ \; --- you cannot isolate \emph{both}. Isolate one, square,
        simplify, and you will find a single radical left. Then square again.
        \vspace{1.10in}

        $x=\blank{1.6cm}$ \quad Check: \blank{4.0cm}
  \item Why does squaring \emph{twice} make checking even more important than usual?
        \writelines{2}
\end{enumerate}

\vspace{4pt}
\textbf{Part 2 --- Prove the rules you have been using.} Assume $A$ and $B$ are real expressions.
\begin{enumerate}[label=(\alph*), leftmargin=1.9em, itemsep=6pt, start=4]
  \item Suppose $\sqrt{A}=B$ has a solution. Explain why $B$ can never be negative there --- and
        then explain why the squared equation $A=B^{2}$ has no way of knowing that. \writelines{3}
  \item Now the odd case. Every real number has exactly \emph{one} real cube root. Use that single
        fact to explain why $A^{3}=B^{3}$ forces $A=B$, and why cubing therefore cannot produce an
        extraneous solution. \writelines{3}
  \item Complete the summary: \emph{raising both sides to an \underline{\hspace{1.6cm}} power can
        add solutions; raising both sides to an \underline{\hspace{1.6cm}} power cannot.}
\end{enumerate}

\vspace{4pt}
\textbf{Part 3 --- Solving backwards, in context.} From Lesson 6.2, a pendulum of length $L$ feet
has period
\[ T(L)=2\pi\sqrt{\tfrac{L}{32}} \ \text{ seconds.} \]
In 6.2 you were given $L$ and found $T$. Now do it the other way --- which means solving a radical
equation.
\begin{enumerate}[label=(\alph*), leftmargin=1.9em, itemsep=6pt, start=7]
  \item A clockmaker wants a pendulum with a period of exactly $2\pi$ seconds. Solve
        $2\pi=2\pi\sqrt{L/32}$ for $L$, and check your answer against 6.2's table.
        \vspace{0.55in}

        $L=\blank{2.0cm}$ feet
  \item Now find the length that gives a period of $\pi$ seconds. \; $L=\blank{2.0cm}$ feet
        \vspace{0.45in}

        The period was cut in half. What happened to the length, and which feature of a square root
        graph explains it? \writelines{2}
  \item An apprentice is asked for a pendulum with period $-\pi$ seconds. There are \emph{two}
        separate reasons to reject that request --- one algebraic, one about the situation. Give
        both, and say why keeping them apart matters. \writelines{3}
\end{enumerate}
\end{tcolorbox}

\end{document}
